\section*{Chapter \ref{ch:fm:managing-a-book-through-drawdown} --- Managing a Book through Drawdown}

\begin{solution}{exo:fm:managing-a-book-through-drawdown:1}
0.8 of a year's volatility at 10\%; 0.27 at 30\%.
\end{solution}

\begin{solution}{exo:fm:managing-a-book-through-drawdown:2}
$A=2\,520$ days: $2\ln2\,520/(1-0)^2=15.7$ years; to $-1$: $15.7/4=3.9$ years.
\end{solution}

\begin{solution}{exo:fm:managing-a-book-through-drawdown:3}
$9.79\times10=98\%$ of healthy strategies over ten years; 66\% of broken ones before their break.
\end{solution}

\begin{solution}{exo:fm:managing-a-book-through-drawdown:4}
Higher: the prior becomes $0.30+0.70/1\,260=0.3006$, and the return favours the broken state by a likelihood ratio of only
$\exp(\tfrac12(S/\sqrt{252})^2)=1.002$, so the posterior is 0.3010. One day's return carries a signal of $S/\sqrt{252}=0.063$ standard deviations.
\end{solution}

\begin{solution}{exo:fm:managing-a-book-through-drawdown:5}
Its expected P\&L is zero after the break, so stopping it forgoes nothing and gains nothing in expectation. It frees the capital and the risk it
uses: the \omterm{def:fm:limits-and-capital-allocation:raroc}{capital charge} while it runs, and the variance it adds to the firm.
\end{solution}

\begin{solution}{exo:fm:managing-a-book-through-drawdown:6}
No: 12\% is below the median ten-year maximum drawdown of a healthy strategy (16.8\%); drawdowns of that depth are normal and say little about a
break.
\end{solution}

\begin{solution}{exo:fm:managing-a-book-through-drawdown:7}
The posterior rule catches broken strategies after a median 1.3 years (false stops 2.24 per hundred years); values for a half-and-half population
are 17.5 (stop-loss), 23.2 (ladder), 31.3 (\omterm{def:fm:managing-a-book-through-drawdown:rules}{time stop}) and 39.3 (posterior) against 10.1 with no rule. When a break turns into losses, every rule
helps and the posterior rule most.
\end{solution}

\begin{solution}{exo:fm:managing-a-book-through-drawdown:8}
The backtest counts the losses avoided and not the recoveries forgone: the stop fires on healthy strategies in ordinary drawdowns (98\% of them
over ten years in the chapter's model). Count false stops and the P\&L lost after them, on strategies known to be healthy, before adopting it.
\end{solution}

\begin{solution}{pb:fm:managing-a-book-through-drawdown:1}
\begin{enumerate}
\item See \cref{def:fm:managing-a-book-through-drawdown:rules}.
\item 16.8\% and 25.4\% of capital.
\item 0.8 and 0.4 of a year's volatility.
\item A healthy strategy's drawdowns scale with its volatility; a threshold in per cent means different things for different strategies.
\item See \cref{prop:fm:managing-a-book-through-drawdown:posterior}.
\item 0.24 for healthy strategies and 0.39 for broken ones (only 13\% of broken paths first reach 8\% after their break).
\item Delay $2\ln A/(S-S_0)^2$ years: 15.7 and 3.9.
\item Capture and fill rates, crowding measures, changes in market structure: observations that move with the edge itself, not with the
strategy's noise.
\item False stops 9.79, 2.09, 6.19 and 2.84 per hundred strategy-years; caught 34\%, 60\%, 80\% and 95\%.
\item 0.72, 2.79, 2.00 and 2.68 years.
\item Live 21.7, 39.5, 54.6, 71.4; broken 15.0, 12.8, 19.5, 18.2; half each 18.3, 26.2, 37.1, 44.8; no rule 79.0, 11.1, 45.1.
\item It stops almost every healthy strategy in an ordinary drawdown, and most broken ones before they break.
\item Stopping a broken strategy at Sharpe 0 saves only its \omterm{def:fm:limits-and-capital-allocation:raroc}{capital charge}; it keeps more when broken strategies lose money (exercise 7).
\item See \cref{def:fm:managing-a-book-through-drawdown:restart}: paper-trade the stopped strategy and restart at reduced size when its paper
P\&L recovers.
\item Reviewing the strategy's thesis and evidence as if new; it uses information other than P\&L.
\item They push managers to gamble for resurrection or to freeze; ladders and \omterm{def:fm:managing-a-book-through-drawdown:restart}{restart rules} soften both.
\item Compare both drawdowns in volatility units and with healthy distributions; look at non-P\&L evidence; cut size only if it points to a
break.
\item The posterior rule with a hazard from the firm's experience, a ladder for size, and a \omterm{def:fm:managing-a-book-through-drawdown:restart}{restart rule}; no depth-only stop.
\item Posterior 0.24 (healthy) and 0.39 (broken) at the 8\% drawdown; 15.7 years to detect a fall from 1 to 0 at one false alarm per ten years.
\item You will be judged on the evidence that your edge is intact, of which your drawdown is a weak part; your size will fall gradually as
evidence against it builds, and a stop will not be the end if the evidence recovers.
\end{enumerate}
\end{solution}

\begin{solution}{iq:fm:managing-a-book-through-drawdown:1}
Not by itself: 15\% is 1.5 years' volatility, within the normal range of a Sharpe-1 strategy's ten-year maximum drawdown (median about 17\%).
\iqlookfor{volatility units and the healthy distribution.}
\end{solution}

\begin{solution}{iq:fm:managing-a-book-through-drawdown:2}
A stop ends trading at one threshold; a ladder cuts size in steps. The ladder has fewer false stops and responds before the stop.
\iqlookfor{graduated response.}
\end{solution}

\begin{solution}{iq:fm:managing-a-book-through-drawdown:3}
KL per day $=(S-S_0)^2/(2\times252)=1/504$; delay $\approx\ln A\times504$ days, about 16 years at $A=2\,520$.
\iqlookfor{information per day and log of the false-alarm period.}
\end{solution}

\begin{solution}{iq:fm:managing-a-book-through-drawdown:4}
Predict $\tilde p=p+(1-p)h$, then Bayes with the two Gaussian likelihoods (\cref{prop:fm:managing-a-book-through-drawdown:posterior}).
\iqlookfor{the prediction and update steps.}
\end{solution}

\begin{solution}{iq:fm:managing-a-book-through-drawdown:5}
How often it would have stopped strategies that went on to recover, and what those recoveries were worth.
\iqlookfor{false stops and forgone recoveries.}
\end{solution}

\begin{solution}{iq:fm:managing-a-book-through-drawdown:6}
Capture, fill and hit rates, signal decay, crowding and flows: they measure the edge directly, with far less noise per day than P\&L.
\iqlookfor{higher signal-to-noise observations.}
\end{solution}
